\documentclass[letterpaper,11pt]{article}

\usepackage[T1]{fontenc}
\usepackage[utf8]{inputenc}
\usepackage[bitstream-charter]{mathdesign}
\usepackage[top=0.17in,bottom=0.22in,left=0.38in,right=0.38in]{geometry}
\usepackage{titlesec}
\usepackage{enumitem}
\usepackage[hidelinks]{hyperref}
\usepackage{fancyhdr}
\usepackage[english]{babel}
\usepackage{tabularx}
\usepackage{array}
\usepackage[protrusion=true,expansion=false]{microtype}
\ifdefined\pdfgentounicode\input{glyphtounicode}\fi

\pagestyle{fancy}
\fancyhf{}
\renewcommand{\headrulewidth}{0pt}
\renewcommand{\footrulewidth}{0pt}
\setlength{\footskip}{12pt}

\urlstyle{same}
\raggedright
\setlength{\tabcolsep}{0in}
\linespread{0.97}
\ifdefined\pdfgentounicode\pdfgentounicode=1\fi

%----------SECTION STYLE: bold serif title + gap + full-width rule below----------
\titleformat{\section}
  {\large\bfseries}
  {}{0em}{}
  [\vspace{2pt}\titlerule]
\titlespacing*{\section}{0pt}{4pt}{2pt}

%----------LIST STYLE: small bullet, hanging indent, target line height----------
\renewcommand\labelitemi{$\vcenter{\hbox{\scriptsize$\bullet$}}$}
\setlist[itemize]{
  leftmargin=0.18in,
  labelsep=0.08in,
  itemsep=0pt,
  topsep=1.5pt,
  parsep=0pt,
  partopsep=0pt
}

%----------ENTRY COMMANDS----------
% Education: bold degree + school left, dates right
\newcommand{\eduItem}[3]{%
  \begin{tabular*}{\textwidth}{@{}l@{\extracolsep{\fill}}r@{}}
    \textbf{#1}, #2 & #3 \\
  \end{tabular*}\par
}

% Experience: Company (bold), Role, | Location (italic) left; Dates right
\newcommand{\expItem}[4]{%
  \vspace{2pt}
  \begin{tabular*}{\textwidth}{@{}l@{\extracolsep{\fill}}r@{}}
    \textbf{#1}, #2 \textit{| #3} & #4 \\
  \end{tabular*}\par\vspace{-3pt}
}

% Project: name (bold) left, tech stack right, same line
\newcommand{\projItem}[2]{%
  \vspace{2pt}
  \begin{tabular*}{\textwidth}{@{}l@{\extracolsep{\fill}}r@{}}
    \textbf{#1} & #2 \\
  \end{tabular*}\par\vspace{-3pt}
}

% Skills: bold label run-in
\newcommand{\skillLine}[2]{\textbf{#1:} #2\par\vspace{1.5pt}}

% Compact inline block with slight left indent (publications, open source & community)
\newcommand{\inlineBlock}[1]{{\setlength{\leftskip}{0.18in}\small #1\par}}

\begin{document}

%----------HEADING----------
\begin{center}
    {\fontsize{18}{20}\selectfont\bfseries DEVA ANAND}\\ \vspace{3pt}
    {\small
    \href{mailto:devaanand@umass.edu}{devaanand@umass.edu}
    \enspace$\bullet$\enspace
    (413) 315-1715
    \enspace$\bullet$\enspace
    \href{https://devaanand.com}{devaanand.com}
    \enspace$\bullet$\enspace
    \href{https://linkedin.com/in/devaaa/}{linkedin.com/in/devaaa}
    \enspace$\bullet$\enspace
    \href{https://github.com/Deva-1903}{github.com/Deva-1903}}
\end{center}
\vspace{-8pt}

%----------EDUCATION----------
\section{Education}
\eduItem{MS, Computer Science}{University of Massachusetts Amherst}{2025--2027}
{\small\textit{Advanced NLP, Advanced Machine Learning, Optimization in CS, Systems for Data Science, Performance Engineering}}\par
\vspace{3pt}
\eduItem{BE, Computer Science and Engineering}{Vels Institute of Science, Technology and Advanced Studies}{2019--2023}

%----------EXPERIENCE----------
\section{Professional Experience}

\expItem{TherAlign Health}{Founding Software Engineer Intern}{Remote}{Jul 2026 -- Present}
\begin{itemize}
  \item Founding engineer on a provider-facing \textbf{SMART on FHIR} clinical decision support app (Next.js + Firebase/GCP) that surfaces covered, guideline-backed medication alternatives in the EHR prescribing flow to prevent prior authorization.
  \item Extending a \textbf{Firebase Cloud Functions} backend that integrates multiple clinical data sources, \textbf{RxNorm} drug normalization, formulary/prior-authorization data, and PubMed/guideline evidence, into structured medication-alternative results, with \textbf{Gemini} constrained to synthesis while the backend deterministically controls eligibility and coverage.
\end{itemize}

\expItem{Culture \& Morality Lab, UMass Amherst}{Lab Assistant}{Amherst, MA}{Jul 2026 -- Present}
\begin{itemize}
  \item Sole engineer building the lab's CCR (Contextualized Construct Representations) text-analysis platform (Python, \textbf{FastAPI}, sentence-transformers, React), turning a published NLP method into a self-serve pipeline that ingests research corpora and produces analysis-ready construct scores.
  \item Designed versioned research-data infrastructure, a YAML \textbf{model registry} as single source of truth, an append-only construct library (\textbf{99 constructs from 38 questionnaires}) with item hashes and immutable per-run snapshots, and a structured \textbf{data-quality} warning system, and grew the backend suite to \textbf{40 hermetic tests} running in CI without ML dependencies.
\end{itemize}

\expItem{\href{https://www.wysa.io}{Wysa}}{Backend Engineer}{Bengaluru, India}{Jul 2023 -- Jul 2025}
\begin{itemize}
  \item Engineered a multi-tenant NHS eTriage backend (Node.js, MongoDB) for \textbf{8+ UK clinical clients}, processing \textbf{10{,}000+ monthly triage submissions} by integrating data across services and Mayden iaptus \textbf{REST} APIs; contributed to \textbf{\$340K+} in revenue.
  \item Architected a full-stack ML training-data annotation platform (React, Node.js, MongoDB) used daily by the AI team, replacing spreadsheet-based labeling with a structured workflow that cut manual annotation by \textbf{100+ hours/month} and lifted labeling throughput by \textbf{60\%}.
\end{itemize}

%----------PROJECTS----------
\section{Projects}

\projItem{Spark ETL Performance Optimization {\normalfont\small(\href{https://github.com/Deva-1903/Spark-ETL-Optimization}{GitHub})}}{PySpark, Spark UI, SQL, NYC TLC Dataset}
\begin{itemize}
  \item Optimized a \textbf{PySpark} ETL + analytics pipeline over a $\sim$1.4B-row NYC TLC dataset with broadcast joins, Adaptive Query Execution, skew-join handling, shuffle-partition tuning, schema-aware column-pruned reads, and year/month partitioning, cutting end-to-end ETL runtime \textbf{$\sim$25\%} (17--19 min to 14.3 min) and the tail-latency ratio \textbf{2.14x to 1.74x}; profiled stage timings in the Spark UI.
\end{itemize}

\projItem{AgenticSearch -- Provenance-First Entity Discovery {\normalfont\small(\href{https://github.com/Deva-1903/ciir_agentic_search}{GitHub} $\vert$ \href{https://agentic-search-negglszkwa-uc.a.run.app}{Demo})}}{Python, FastAPI, OpenAI / Groq, SQLite}
\begin{itemize}
  \item Built a multi-stage Python/\textbf{FastAPI} data pipeline that turns open-ended queries into \textbf{structured entity tables} via typed retrieval planning, page rerank, and deterministic extraction with LLM fallback, returning \textbf{cell-level provenance} (value, source, snippet, confidence) for every value; hardened with \textbf{150+ automated tests}.
\end{itemize}

\projItem{KG2RAG-Enhanced -- Multi-Hop QA on HotpotQA {\normalfont\small(\href{https://github.com/Deva-1903/KG2RAG-685-NLP}{GitHub})}}{Python, Ollama / LLaMA-3, sentence-transformers, RRF}
\begin{itemize}
  \item Built the multi-view retrieval component of a KG-guided \textbf{RAG} data pipeline (query-side decomposition, RRF fusion, cross-encoder reranking, MMR) on HotpotQA, lifting supporting-fact recall \textbf{+2.68\%} (54.53 to 57.21) over the KG\textsuperscript{2}RAG baseline on a 4{,}905-question evaluation.
\end{itemize}

%----------PUBLICATIONS----------
\section{Publications}
\vspace{2pt}
\inlineBlock{- \textit{``Alzheimer's Disease Classification using Transfer Learning,''} IEEE CONIT 2023, first author, deep transfer learning on neuroimaging data. {\normalfont\small(\href{https://ieeexplore.ieee.org/document/10205760}{Paper})}}

%----------SKILLS----------
\section{Skills}
\skillLine{Languages \& Databases}{Python, SQL, TypeScript, JavaScript, Java, C++, Bash, PostgreSQL, MongoDB, SQLite}
\skillLine{Data \& Pipelines}{PySpark, Spark, ETL, distributed data processing, data modeling, partitioning, data-quality checks, Hadoop / MapReduce}
\skillLine{Backend \& APIs}{FastAPI, Node.js, Express, Next.js, REST APIs, async services, pytest / integration testing}
\skillLine{ML \& AI}{PyTorch, sentence-transformers, Hugging Face Transformers, RAG, LLM agents, evaluation}
\skillLine{Cloud \& DevOps}{Azure, AWS, GCP, Docker, Git, GitHub Actions, CI/CD, Linux}

\end{document}
